\chapter{Appendix: Canonical neighborhoods}
\label{chap:canonical-neighborhoods}
\label{sect:canonnbhd}


Recall that an $\epsilon$-neck structure on a Riemannian manifold
$(N,g)$ centered at a point $x\in N$ is a diffeomorphism $\psi\colon
S^2\times (-\epsilon^{-1},\epsilon^{-1})\to N$ with the property
that $x\in \psi(S^2\times\{0\})$ and the property that $R(x)\psi^*g$
is within $\epsilon$ in the $C^{[1/\epsilon]}$-topology of the
product metric $h_0\times ds^2$, where $h_0$ is the round metric on
$S^2$ of scalar curvature $1$ and $ds^2$ is the Euclidean metric on
the interval. Recall that the {\em scale} of the $\epsilon$-neck is
$R(x)^{-1/2}$. We define $s=s_N\colon N\to
(-\epsilon^{-1},\epsilon^{-1})$ as the composition of $\psi^{-1}$
followed by the projection to the second factor.


\section{Shortening curves}

\begin{lemma}[Shortening a curve through the negative end of a neck]
\label{lem:neck-shortening-curve}
\uses{def:epsilon-neck}
\label{curveshort}

The following holds for all $\epsilon>0$ sufficiently small. Suppose
that $(M,g)$ is a Riemannian manifold and that $N\subset M$ is an
$\epsilon$-neck centered at $x$. Let $S(x)$ be the central
two-sphere of this neck and suppose that $S(x)$ separates $M$. Let
$y\in M$. Orient $s$ so that $y$ lies in the closure of the positive
side of $S(x)$. Let $\gamma\colon [0,a]\to M$ be a rectifiable curve
from $x$ to $y$. If $\gamma$ contains a point of
$s^{-1}(-\epsilon^{-1},-\epsilon^{-1}/2)$ then there is a
rectifiable curve from $x$ to $y$ contained in the closure of the
positive side of $S(x)$ whose length is at most the length of
$\gamma$ minus $\frac{1}{2}\epsilon^{-1}R(x)^{-1/2}$.
\end{lemma}


\begin{proof}
\sketch
\uses{def:epsilon-neck}
\end{proof}



\section{The geometry of an $\epsilon$-neck}


\begin{lemma}[Curvature pinching in an $\epsilon$-neck]
\label{lem:epsilon-neck-curvature-pinching}
\uses{def:epsilon-neck}
\label{neckcurv}

For any $0<\alpha<1/8$  there is $\epsilon_1=\epsilon_1(\alpha)>0$ such that
the following two conditions hold for all $0<\epsilon\le \epsilon_1$.
\begin{enumerate} \item[(1)] If $(N,g)$ is an $\epsilon$-neck centered at $x$ of
scale one (i.e., with $R(x)=1$) then the principal sectional
curvatures at any point of $N$ are within $\alpha/6$ of
$\{1/2,0,0\}$. In particular, for any $y\in N$ we have
$$(1-\alpha)\le R(y)\le (1+\alpha).$$
\item[(2)] There is unique two-plane of maximal sectional curvature at
every point of an $\epsilon$-neck, and the angle between the
distribution of two-planes of maximal sectional curvature and the
two-plane field tangent to the family of two-spheres of the
$\epsilon$-neck structure is everywhere less than $\alpha$.
\end{enumerate}
\end{lemma}

\begin{proof}
\sketch
\end{proof}


\begin{corollary}[An almost-round embedded sphere is isotopic to a neck fiber]
\label{cor:embedded-sphere-isotopic-to-neck-fiber}
\uses{def:epsilon-neck}
\label{s2isotopic}

The following holds for any $\epsilon>0$ sufficiently small. Suppose
that $(N,g)$ is an $\epsilon$-neck and we have  and an embedding
$f\colon S^2 \to N$ with the property that the restriction of $g$ to
the image of this embedding is within $\epsilon$ in the
$C^{[1/\epsilon]}$-topology to the round metric $h_0$ of scalar
curvature one on $S^2$ and with the norm of the second fundamental
form less than $\epsilon$. Then the two-sphere $f(S^2)$ is isotopic
in $N$ to any member of the family of two-spheres coming from the
$\epsilon$-neck structure on $N$.
\end{corollary}



\begin{proof}
\sketch
\uses{lem:epsilon-neck-curvature-pinching}
\end{proof}


\begin{lemma}[Minimal geodesics in a neck nearly align with the axis direction]
\label{lem:neck-geodesic-directions}
\uses{def:epsilon-neck}
\label{directions}

 For any $\alpha>0$  there is $\epsilon_2=\epsilon_2(\alpha)>0$
such that the following hold for all $0<\epsilon\le \epsilon_2$.
Suppose that $(N,g)$ is an $\epsilon$-neck centered at $x$ and
$R(x)=1$. Suppose that $\gamma$ is a minimal geodesic in  $N$ from
$p$ to $q$. We suppose that $\gamma$ is parameterized by arc length,
is of length $\ell>\epsilon^{-1}/100$, and that $s(p)<s(q)$. Then
for all $s$ in the domain of definition of $\gamma$ we have
$$|\gamma'(s)-(\partial/\partial s)|_g<\alpha.$$
 In particular, the angle between $\gamma'$ and $\partial/\partial
s$ is less than $2\alpha$.
 Also,  any
member $S^2$ of the family of two-spheres in the $N$ has intrinsic
diameter at most $(1+\alpha)\sqrt{2}\pi$.
\end{lemma}

\begin{proof}
\sketch
\end{proof}


\begin{corollary}[Distance and diameter estimates in an $\epsilon$-neck]
\label{cor:neck-distance-estimates}
\uses{def:epsilon-neck}
\label{neckgeo}

For any $\alpha>0$ there is $\epsilon_3=\epsilon_3(\alpha)>0$ such
that the  following hold for any $0<\epsilon\le \epsilon_3$ and any
$\epsilon$-neck $N$ of scale $1$ centered at $x$.
\begin{enumerate}
\item[(1)] Suppose that $p$ and $q$ are points of  $N$
with either $|s(q)-s(p)|\ge \epsilon^{-1}/100$ or $d(p,q)\ge
\epsilon^{-1}/100$. Then  we have
$$(1-\alpha)|s(q)-s(p)|\le d(p,q)\le (1+\alpha)|s(q)-s(p)|.$$
\item[(2)]
$$B(x,(1-\alpha)\epsilon^{-1})\subset
N\subset B(x,(1+\alpha)\epsilon^{-1}).$$
\item[(3)] Any geodesic that exits from both ends of $N$ has length at least $2(1-\alpha)
\epsilon^{-1}$.\end{enumerate}
\end{corollary}



\begin{corollary}[A long shortest geodesic crosses each neck sphere at most once]
\label{cor:shortest-geodesic-crosses-neck-once}
\uses{def:epsilon-neck}
\label{compint}

The following holds for all $\epsilon>0$ sufficiently small. Let $N$
be an $\epsilon$-neck centered at $x$. If $\gamma$ is a shortest
geodesic in $N$ between its endpoints and if
$|\gamma|>R(x)^{-1/2}\epsilon^{-1}/100$, then $\gamma$ crosses each
two-sphere in the neck structure on $N$ at most once.
\end{corollary}






There is a closely related lemma.

\begin{lemma}[A shortest geodesic crossing a neck at most once, in general]
\label{lem:geodesic-crosses-neck-once-general}
\uses{def:epsilon-neck}
\label{genint}

The following holds for every $\epsilon>0$ sufficiently small.
Suppose that $(M,g)$ is a Riemannian manifold and that $N\subset M$
is an $\epsilon$-neck centered at $x$ and suppose that $\gamma$ is a
shortest geodesic in $M$ between its endpoints and that the length
of every component of $N\cap|\gamma|$ has length at least $
R(x)^{-1/2}\epsilon^{-1}/8$. Then $\gamma$ crosses each two-sphere
in the neck structure on $N$ at most once.
\end{lemma}


\begin{proof}
\sketch
\uses{cor:shortest-geodesic-crosses-neck-once,cor:neck-distance-estimates,lem:neck-geodesic-directions}
\end{proof}


\begin{corollary}[A minimal geodesic through the core of a cap]
\label{cor:minimal-geodesic-through-core}
\uses{def:epsilon-neck, def:c-epsilon-cap}
\label{coreint}

The following holds for all $\epsilon>0$ sufficiently small and any
$C<\infty$. Let $X$ be an $(C,\epsilon)$-cap in a complete
Riemannian manifold $(M,g)$, and let $Y$ be its core and let $S$ be
the central two-sphere of the $\epsilon$-neck $N=X-\overline Y$. We
orient the $s$-direction in $N$ so that $Y$ lies off the negative
end of $N$. Let $\widehat Y$ be the union of $\overline Y$ and the
closed negative half of $N$ and let $S$ be the boundary of $\widehat
Y$. Suppose that $\gamma$ is a minimal geodesic in $(M,g)$ that
contains a point of the core $Y$. Then the intersection of $\gamma$
with $\widehat Y$ is an interval containing an endpoint of $\gamma$.
\end{corollary}



\begin{proof}
\sketch
\uses{cor:neck-distance-estimates,cor:shortest-geodesic-crosses-neck-once}
\end{proof}




We also wish to compare distances from points outside the neck with
distances in the neck.

\begin{lemma}[Distance comparison from a point outside the neck]
\label{lem:neck-distance-comparison-external-point}
\uses{def:epsilon-neck}
\label{neckdistcomp}

Given $0<\alpha<1$ there is $\epsilon_4=\epsilon_4(\alpha)>0$ such
that the following holds for any $0<\epsilon\le \epsilon_4$. Suppose
that $N$ is an $\epsilon$-neck centered at $x$ in a connected
manifold $M$ (here we are not assuming that $R(x)=1$). We suppose
that the central $2$-sphere of $N$ separates $M$. Let $z$ be a point
outside of  the middle two-thirds of $N$ and lying on the negative
side of the central $2$-sphere of $N$. (We allow both the case when
$z\in N$ and when $z\not\in N$.) Let $p$ be a point in the middle
half of $N$. Let $\mu\colon [0,a]\to N$ be a straight line segment
(with respect to the standard product metric) in the positive
$s$-direction in $N$ beginning at $p$ and ending at a point $q$ of
$N$. Then
$$(1-\alpha)(s(q)-s(p))\le d(z,q)-d(z,p)\le (1+\alpha)(s(q)-s(p)).$$
\end{lemma}


\begin{proof}
\sketch
\end{proof}



\noindent{\bf N.B.} It is important that the central two-sphere of
$N$ separates the ambient manifold $M$. Otherwise, there may be
shorter geodesics from $z$ to $q$ entering the other end of $N$.

\begin{lemma}[Metric-ball boundary meets a neck in a topological sphere]
\label{lem:metric-sphere-meets-neck-topologically}
\uses{def:epsilon-neck}
\label{topiso}

Given any $\alpha>0$  there is $\epsilon(\alpha)>0$ such that  the
following holds for any $0<\epsilon\le \epsilon(\alpha)$. Suppose
that $N$ is an $\epsilon$-neck centered at $x$ in a connected
manifold $M$ (here we are not assuming that $R(x)=1$) and that $z$
is a point outside the middle two-thirds of $N$. We suppose that the
central two-sphere of $N$ separates $M$. Let $p$ be a point in the
middle sixth of $N$ at distance $d$ from $z$. Then the intersection
of the boundary of the metric ball $B(z,d)$ with $N$ is a
topological $2$-sphere contained in the middle quarter of $N$ that
maps homeomorphically onto $S^2$ under the projection mapping $N\to
S^2$ determined by the $\epsilon$-neck structure. Furthermore, if
$p'\in \partial B(z,d)$ then $|s(p)-s(p')|<\alpha
R(x)^{-1/2}\epsilon^{-1}$.
\end{lemma}

\begin{proof}
\sketch
\uses{lem:neck-distance-comparison-external-point}
\end{proof}






\section{Overlapping $\epsilon$-necks}

The subject of this section is the internal geometric properties of
$\epsilon$-necks and of intersections of $\epsilon$-necks. We assume that
$\epsilon\le 1/200$.


\begin{proposition}[Properties of overlapping $\epsilon$-necks]
\label{prop:overlapping-necks-properties}
\uses{def:epsilon-neck}
\label{S2intersection}

Given $0<\alpha\le 10^{-2}$, there is
$\epsilon_5=\epsilon_5(\alpha)>0$ such that the following hold for
all $0<\epsilon\le \epsilon_5$. Let $N$ and $N'$ be $\epsilon$-necks
centered at $x$ and $x'$, respectively, in a Riemannian manifold
$X$:
\begin{enumerate}
\item[(1)] If $N\cap N'\not=\emptyset$ then
$1-\alpha<R(x)/R(x')<1+\alpha$. In particular, denoting the scales
of $N$ and $N'$ by $h$ and $h'$ we have
$$1-\alpha<\frac{h}{h'}<1+\alpha.$$
\item[(2)] Suppose $y\in N\cap N'$ and $S$ and $S'$ are the two-spheres in the $\epsilon$-neck
structures on $N$ and $N'$, respectively, passing through $y$. Then
the angle between $TS_y$ and $TS'_y$ is less than $\alpha$.
\item[(3)] Suppose that $y\in N\cap N'$. Denote by $\partial/\partial s_N$ and $\partial/\partial s_{N'}$
the tangent vectors in the $\epsilon$-neck structures of $N$ and
$N'$, respectively. Then at the point $y$, either
$$|R(x)^{1/2}(\partial /\partial s_N)-R(x')^{1/2}\partial/\partial s_{N'}|<\alpha$$
or
$$|R(x)^{1/2}(\partial /\partial s_N)+R(x')^{1/2}\partial/\partial s_{N'}|<\alpha.$$
 \item[(4)] Suppose that
one of the two-spheres $S'$ of the $\epsilon$-neck structure on $N'$
is completely contained in $N$. Then $S'$ is a section of the
projection mapping on the first factor
$$p_1\colon S^2\times(-\epsilon^{-1},\epsilon^{-1})\to S^2.$$
In particular, $S'$ is isotopic in $N$ to any one of the two-spheres
of the $\epsilon$-neck structure on $N$ by an isotopy that moves all
points in the interval directions.
\item[(5)] If $N\cap N'$ contains a point $y$ with $(-0.9)\epsilon^{-1}\le s_N(y)\le
(0.9)\epsilon^{-1}$, then there is a point $y'\in N\cap N'$ such
that
$$-(0.96)\epsilon^{-1}\le s_N(y')\le (0.96)\epsilon^{-1}$$
$$-(0.96)\epsilon^{-1}\le s_{N'}(y')\le (0.96)\epsilon^{-1}.$$
The two-sphere $S(y')$ in the neck structure on $N$ through $y'$ is contained
in $N'$ and the two-sphere $S'(y')$ in the neck structure on $N'$ through $y'$
is contained in $N$. Furthermore, $S(y')$ and $S'(y')$ are isotopic in $N\cap
N'$. Lastly, $N\cap N'$ is diffeomorphic to $S^2\times (0,1)$ under a
diffeomorphism mapping $S(y)$ to $S^2\times \{1/2\}$.
\end{enumerate}
\end{proposition}


\begin{proof}
\sketch
\uses{cor:embedded-sphere-isotopic-to-neck-fiber,lem:epsilon-neck-curvature-pinching,lem:neck-geodesic-directions,cor:neck-distance-estimates}
\end{proof}






\section{Regions covered by $\epsilon$-necks and
$(C,\epsilon)$-caps}


Here we fix $0<\epsilon\le 1/200$ sufficiently small so that all the
results in the previous two sections hold with $\alpha=10^{-2}$.

\subsection{Chains of $\epsilon$-necks}

\begin{definition}[Finite and infinite chains of $\epsilon$-necks]
\label{def:chain-of-epsilon-necks}
\uses{def:epsilon-neck}

Let $(X,g)$ be a Riemannian manifold. By a {\em finite chain} of
$\epsilon$-necks in $(X,g)$, we mean a sequence $N_a,\ldots,N_{b}$,
of $\epsilon$-necks in $(X,g)$ such that: \begin{enumerate} \item
for all $i,\ a\le i<b$, the intersection $N_i\cap N_{i+1}$ contains
the positive-most quarter of $N_i$ and the negative-most quarter of
$N_{i+1}$ and is contained in the positive-most three-quarters of
$N_i$ and the negative-most three-quarters of $N_{i+1}$, and
\item  for all $i,\ a<i\le b$, $N_i$ is disjoint from the negative end of
$N_a$.
\end{enumerate}
By an {\em infinite chain} of $\epsilon$-necks in $X$ we mean a
collection $\{N_i\}_{i\in I}$ for some interval $I\subset \Zee$,
infinite in at least one direction, so that for each finite
subinterval $J$ of $I$ the subset of $\{N_i\}_{i\in J}$ is a chain
of $\epsilon$-necks.
\end{definition}




Notice that in an $\epsilon$-chain $N_i\cap N_j=\emptyset$ if
$|i-j|\ge 5$.


\begin{lemma}[Union of a chain of necks is diffeomorphic to $S^2\times(0,1)$]
\label{lem:chain-union-diffeomorphic-cylinder}
\uses{def:chain-of-epsilon-necks}
\label{S2timesI}

The union $U$ of the $N_i$ in a finite or infinite chain of
$\epsilon$-necks is diffeomorphic to $S^2\times (0,1)$. In
particular, it is an $\epsilon$-tube.
\end{lemma}

\begin{proof}
\sketch
\uses{prop:overlapping-necks-properties}
\end{proof}

Notice that the frontier of the union of the necks in a finite
chain, $U=\cup_{a\le i\le b}N_i$, in $M$ is equal to the frontier of
the positive end of $N_b$ union the frontier of the negative end of
$N_a$. Thus, we have:

\begin{corollary}[Frontier of a finite chain of necks]
\label{cor:chain-frontier}
\uses{def:chain-of-epsilon-necks}
\label{chainfrontier}

Let $\{N_a,\ldots,N_b\}$ be a chain of $\epsilon$-necks. If a
connected set $Y$ meets both $U=\cup_{a\le i\le b}N_i$ and its
complement, then $Y$ either contains  points of the frontier of the
negative end $N_a$ or of the positive  end of $N_b$.
\end{corollary}


The next result shows there is no frontier at an infinite end.

\begin{lemma}[Frontier of an infinite chain of necks]
\label{lem:infinite-chain-frontier}
\uses{def:chain-of-epsilon-necks}
\label{inffrontier}

Suppose that $\{N_0,\cdots\}$ is an infinite chain of
$\epsilon$-necks in $M$. Then the frontier of $U=\cup_{i=0}^\infty
N_i$ is the frontier of the negative end of $N_0$.
\end{lemma}

\begin{proof}
\sketch
\uses{cor:neck-distance-estimates}
\end{proof}


 In fact, there is a geometric version of Lemma~\ref{lem:chain-union-diffeomorphic-cylinder}.

\begin{lemma}[Geometric fibration of a neck chain by two-spheres]
\label{lem:chain-fibration-by-spheres}
\uses{def:chain-of-epsilon-necks}

There is $\epsilon_0>0$ such that  the following holds for all
$0<\epsilon\le \epsilon_0$. Suppose that $\{N_j\}_{j\in J}$ is a
chain of $\epsilon$-necks in a Riemannian manifold $M$. Let
$U=\cup_{j\in J}N_j$. Then there exist an interval $I$ and a smooth
map $p\colon U\to I$ such that every fiber of $p$ is a two-sphere,
and if $y$ is in the middle $7/8$'s of $N_j$ then the fiber
$p^{-1}(p(y))$ makes a small angle at every point with the family of
two-spheres in the $\epsilon$-neck $N_j$.
\end{lemma}


\begin{proof}
\sketch
\uses{lem:epsilon-neck-curvature-pinching}
\end{proof}

A finite or infinite chain $\{N_j\}_{j\in J}$ of $\epsilon$-necks is
{\em balanced} provided that for every $j\in J$, not the largest
element of $J$, we have
\begin{equation}\label{distineq}
(0.99)R(x_j)^{-1/2}\epsilon^{-1}\le d(x_j,x_{j+1})\le
(1.01)R(x_j)^{-1/2}\epsilon^{-1},
\end{equation}
where, for each $j$, $x_j$ is the central point of $N_j$.


Notice that in a balanced chain $N_j\cap N_{j'}=\emptyset$ if
$|j-j'|\ge 3$.










\begin{lemma}[Two adjacent necks form a balanced chain]
\label{lem:two-necks-form-balanced-chain}
\uses{def:chain-of-epsilon-necks}
\label{firstchain}

There exists $\epsilon_0>0$ such that for all $0<\epsilon\le
\epsilon_0$ the following is true.  Suppose that $N$ and $N'$ are
$\epsilon$-necks centered at $x$ and $x'$, respectively, in a
Riemannian manifold $M$. Suppose that $x'$ is not contained in $N$
but is contained in the closure of $N$ in $M$. Suppose also that the
two-spheres of the neck structure on $N$ and $N'$ separate $M$.
Then, possibly after reversing the $\epsilon$-neck structures on $N$
and/or $N'$, the pair $\{N,N'\}$ forms a balanced chain.
\end{lemma}

\begin{proof}
\sketch
\uses{cor:neck-distance-estimates}
\end{proof}




\begin{lemma}[Extending a balanced chain by one neck]
\label{lem:extend-balanced-chain}
\uses{def:chain-of-epsilon-necks}
\label{addaneck}

There exists $\epsilon_0>0$ such that for all $0<\epsilon\le
\epsilon_0$ the following is true.  Suppose that
$\{N_a,\ldots,N_b\}$ is a balanced chain in a Riemannian manifold
$M$ with $U=\cup_{i=a}^bN_i$. Suppose that the two-spheres of the
neck structure of $N_a$ separate $M$.
 Suppose that $x$
is a point of the frontier of $U$ contained in the closure of the
plus end of $N_b$ that is also the center of an $\epsilon$-neck $N$.
Then possibly after reversing the direction of $N$, we have that
$\{N_a,\ldots,N_b,N\}$ is a balanced chain. Similarly, if $x$ is in
the closure of the minus end of $N_a$, then (again after possibly
reversing the direction of $N$) we have that $\{N,N_a,\ldots,N_b\}$
is a balanced $\epsilon$-chain.
\end{lemma}

\begin{proof}
\sketch
\uses{lem:two-necks-form-balanced-chain}
\end{proof}

\begin{proposition}[A connected set of neck centers is covered by a balanced chain]
\label{prop:balanced-chain-covers-connected-set}
\uses{def:chain-of-epsilon-necks}
\label{epschains}

There exists $\epsilon_0>0$ such that for all $0<\epsilon\le
\epsilon_0$ the following is true. Let $X$ be a connected subset of
a Riemannian manifold $M$ with the property that every point $x\in
X$ is the center of an $\epsilon$-neck $N(x)$ in $M$. Suppose that
the central two-spheres of these necks do not separate $M$. Then
there is a subset $\{x_i\}$ of $X$ such that the necks $N(x_i)$
(possibly after reversing their $s$-directions) form a balanced
chain of $\epsilon$-necks $\{N(x_i)\}$ whose union $U$ contains $X$.
The union $U$ is diffeomorphic to $S^2\times (0,1)$. It is an
$\epsilon$-tube.
\end{proposition}

\begin{proof}
\sketch
\uses{lem:extend-balanced-chain,lem:infinite-chain-frontier,lem:chain-union-diffeomorphic-cylinder}
\end{proof}

\begin{lemma}[A manifold covered by $\epsilon$-necks is a tube or an $S^2$-fibration over $S^1$]
\label{lem:manifold-covered-by-necks-classification}
\uses{def:epsilon-neck}
\label{S2times01}

The following holds for every $\epsilon>0$ sufficiently small. Let
$(M,g)$ be a connected Riemannian manifold. Suppose that every point
of $M$ is the center of an $\epsilon$-neck. Then either $M$ is
diffeomorphic to $S^2\times (0,1)$ and is an $\epsilon$-tube, or $M$
is diffeomorphic to an $S^2$-fibration over $S^1$.
\end{lemma}

\begin{proof}
\sketch
\uses{prop:balanced-chain-covers-connected-set}
\end{proof}


\section{Subsets of the union of cores of $(C,\epsilon)$-caps
and $\epsilon$-necks.} In this section we fix $0<\epsilon\le 1/200$
so that all the results of this section hold with $\alpha=0.01$.

\begin{proposition}[Connected sets covered by caps and necks]
\label{prop:connected-set-in-caps-and-necks}
\uses{def:epsilon-neck, def:c-epsilon-cap}
\label{Xcontainedin}

For any $C<\infty$ the following holds. Suppose that $X$ is a
connected subset of a Riemannian three-manifold $(M,g)$. Suppose
that every point of $X$ is either the center of an $\epsilon$-neck
or is contained in the core of a $(C,\epsilon)$-cap. Then one of the
following holds:
\begin{enumerate}
\item[(1)] $X$ is contained in a component of $M$ that is the union of
two $(C,\epsilon)$-caps. This component is diffeomorphic to $S^3$,
$\Ar P^3$ or $\Ar P^3\#\Ar P^3$.
\item[(2)] $X$ is contained in a component of $M$ that is a double $C$-capped $\epsilon$-tube.
 This component is diffeomorphic to $S^3$, $\Ar P^3$ or $\Ar
P^3\#\Ar P^3$.
\item[(3)] $X$ is contained in a single $(C,\epsilon)$-cap.
\item[(4)] $X$ is contained in a $C$-capped $\epsilon$-tube.
\item[(5)] $X$ is contained in an $\epsilon$-tube.
\item[(6)] $X$ is contained in a component of $M$ that is an
$\epsilon$-fibration, which itself is a union of $\epsilon$-necks.
\end{enumerate}
\end{proposition}


\begin{proof}
\sketch
\uses{prop:balanced-chain-covers-connected-set,lem:manifold-covered-by-necks-classification,prop:overlapping-necks-properties,cor:neck-distance-estimates,lem:epsilon-neck-curvature-pinching,lem:infinite-chain-frontier,def:c-epsilon-cap}
We divide the proof into two cases: Case I: There is a point of $X$
contained in the core of a $(C,\epsilon)$-cap. Case II: Every point
of $X$ is the center of an $\epsilon$-neck.

\noindent{\bf Case I:} We begin the study of this case with a claim.

Now let us turn to the proof of the proposition. We suppose first
that there is a point $x_0\in X$ that is contained in the core of a
$(C,\epsilon)$-cap. Applying the previous claim, we can find a
$(C,\epsilon)$-cap $C_0$ containing $x_0$ with the property that no
point of $X$ contained in the  frontier of $C_0$ is contained in the
closure of the core of a $(C,\epsilon)$-cap $C_1$ that contains
$C_0$.

There are three possibilities to examine:
\begin{enumerate}
\item[(i)] $X$ is disjoint from the frontier of $C_0$.
\item[(ii)] $X$ meets the frontier of $C_0$ but every point of this
intersection is the center of an $\epsilon$-neck.
\item[(iii)] There is a point of the intersection of $X$
with the frontier of $C_0$ that is contained in the core of
$(C,\epsilon)$-cap.
\end{enumerate}

In the first case, since $X$ is connected, it is contained in $C_0$.
In the second case we let $N_1$ be an $\epsilon$-neck centered at a
point of the intersection of $X$ with the frontier of $C_0$, and we
replace $C_0$ by $C_0\cup N_1$ and repeat the argument at the
frontier of $C_0\cup N_1$. We continue in this way creating $C_0$
union a balanced chain of $\epsilon$-necks $C_0\cup N_1\cup N_2\cup
\cdots\cup N_k$. At each step it is possible that either there is no
point of the frontier containing a point of $X$, in which case the
union, which is a $C$-capped $\epsilon$-tube, contains $X$. Another
possibility is that we can repeat the process forever creating a
$C$-capped infinite $\epsilon$-tube. By Lemma~\ref{lem:infinite-chain-frontier} this
union is a component of $M$ and hence contains $X$.

We have shown that one of following holds:
\begin{enumerate}
\item[(a)] There is a $(C,\epsilon)$-cap that contains $X$.
\item[(b)] There is a finite or infinite $C$-capped $\epsilon$-tube that
contains $X$.
\item[(c)] There is a $(C,\epsilon)$-cap or a finite $C$-capped
$\epsilon$-tube $\widetilde C$ containing a point of $X$ and there
is a point of the intersection of $X$ with the frontier of
$\widetilde C$ that is contained in the core of a
$(C,\epsilon)$-cap.
\end{enumerate}

In the first two cases we have established the proposition. Let us
examine the third case in more detail. Let $N_0\subset C_0$ be the
$\epsilon$-neck that is the complement of the closure of the core of
$C_0$.
 First notice that by
Lemma~\ref{lem:manifold-covered-by-necks-classification} the union $N_0\cup N_1\cup\cdots\cup N_k$ is
diffeomorphic to $S^2\times (0,1)$, with the two-spheres coming from
the $\epsilon$-neck structure of each $N_i$ being isotopic to the
two-sphere factor in this product structure. It follows immediately
that $\widetilde C$ is diffeomorphic to $C_0$. Let $C'$ be a
$(C,\epsilon)$-cap whose core contains a point of the intersection
of $X$ with the frontier of $\widetilde C$. We use the terminology
`the core of $\widetilde C$' to mean $\widetilde C\setminus N_k$.
Notice that if $k=0$, this is exactly the core of $C_0$. To complete
the proof of the result we must show that the following hold:

The last thing to do in this case in order to prove the proposition
in Case I is to show that $\widetilde C\cup C'$ is diffeomorphic to
$S^3$, $\Ar P^3$, or $\Ar P^3\#\Ar P^3$. The reason for this is that
$\widetilde C$ is diffeomorphic to $C_0$;  hence $\widetilde C$
either is diffeomorphic to an open three-ball or to a punctured $\Ar
P^3$. Thus, the frontier of $C'$ in $\widetilde C$ is a two-sphere
that bounds either a compact three-ball or the complement of an open
three-ball in $\Ar P^3$. Since $C'$ itself is diffeomorphic either
to a three-ball or to a punctured $\Ar P^3$, the result follows.

\noindent{\bf Case II:} Suppose that every point of $X$ is the
center of an $\epsilon$-neck. Then if the two-spheres of these necks
separate $M$, it follows from Proposition~\ref{prop:balanced-chain-covers-connected-set} that $X$
is contained in an $\epsilon$-tube in $M$.

It remains to consider the case when the two-spheres of these necks
do not separate $M$. As in the case when the two-spheres separate,
we begin building a balanced chain $\epsilon$-necks with each neck
in the chain centered at a point of $X$. Either this construction
terminates after a finite number of steps in a finite
$\epsilon$-chain whose union contains $X$, or it can be continued
infinitely often creating an infinite $\epsilon$ chain containing
$X$ or at some finite stage (possibly after reversing the indexing
and the $s$-directions of the necks)  we have a balanced
$\epsilon$-chain $N_a\cup\cdots\cup N_{b-1}$ and a point of the
intersection of $X$ with the frontier of the positive end of
$N_{b-1}$ that is the center of an $\epsilon$-neck $N_b$ with the
property that $N_b$ meets the negative end of $N_a$. Intuitively,
the chain wraps around on itself like a snake eating its tail. If
the intersection of $N_a\cap N_b$ contains a point $x$ with
$s_{N_a}(x)\ge -(0.9)\epsilon^{-1}$, then according to
Proposition~\ref{prop:overlapping-necks-properties} the intersection of $N_a\cap N_b$
is diffeomorphic to $S^2\times (0,1)$ and the two-sphere in this
product structure is isotopic in $N_a$ to the central two-sphere of
$N_a$ and is isotopic in $N_b$ to the central two-sphere of $N_b$.
In this case it is clear that $N_a\cup \cdots \cup N_b$ is a
component of $M$ that is an $\epsilon$-fibration.

We examine the possibility that the intersection $N_a\cap N_b$
contains some points in the negative end of $N_a$ but is  contained
in $s_{N_a}^{-1}((-\epsilon^{-1},-(0.9)\epsilon^{-1}))$. Set $A=
s_{N_a}^{-1}((-\epsilon^{-1},-(0.8)\epsilon^{-1}))$. Notice that
since $X$ is connected and $X$ contains a point in the frontier of
the positive end of $N_a$ (since we have added at least one neck at
this end), it follows that $X$ contains points in $s_{N_a}^{-1}(s)$
for all $s\in [0,\epsilon^{-1})$.
 If there are
no points of $X$ in  $A$, then we replace $N_a$ by an
$\epsilon$-neck $N_a'$ centered at a point of
$s_{N_a}^{-1}\left((0.15)\epsilon^{-1}\right)\cap X$. Clearly, by
Lemma~\ref{lem:manifold-covered-by-necks-classification} $N'_a$ contains
$s_{N_a}^{-1}(-(0.8)\epsilon{-1},\epsilon^{-1})$ and is disjoint
from $s_{N_a}^{-1}((-\epsilon^{-1},-(0.9)\epsilon^{-1})$, so that
$N_a',N_{a+1}, \ldots, N_b$ is a chain of $\epsilon$-necks
containing $X$. If there is a point of $X\cap A$, then we let
$N_{b+1}$ be a neck centered at this point. Clearly,
$N_a\cup\cdots\cup N_{b+1}$ is a component, $M_0$, of $M$ containing
$X$. The preimage in the universal covering of $M_0$ is a chain of
$\epsilon$-necks infinite in both directions. That is to say, the
universal covering of $M_0$ is an $\epsilon$-tube. Furthermore, each
point in the universal cover of $M_0$ is the center of an
$\epsilon$-neck that is disjoint from all its non-trivial covering
translates. Hence, the quotient $M_0$ is an $\epsilon$-fibration.

We have now completed the proof of Proposition~\ref{prop:connected-set-in-caps-and-necks}.
\end{proof}


\begin{claim}[No infinite nested chain of $(C,\epsilon)$-caps]
\label{claim:no-infinite-chain-of-caps}
It is not possible to have an infinite chain of $(C,\epsilon)$-caps
$C_0\subset C_1\subset \cdots$ in $M$ with the property that for
each $i\ge 1$, the closure of the core of $C_i$ contains a point of
the frontier of $C_{i-1}$
\end{claim}

\begin{proof}
\sketch
\uses{lem:epsilon-neck-curvature-pinching,cor:neck-distance-estimates,prop:overlapping-necks-properties,def:c-epsilon-cap}
We argue by contradiction. Suppose there is such an infinite chain.
Fix a point $x_0\in C_0$ and let $Q_0=R(x_0)$. For each $i\ge 1$ let
$x_i$ be a point in the frontier of $C_{i-1}$ that is contained in
the closure of the core of $C_i$. For each $i$ let $N_i$ be the
$\epsilon$-neck in $C_{i}$ that is the complement of the closure of
its core. We orient the $s_{N_{i}}$-direction so that the core of
$C_{i}$ lies off the negative end of $N_{i}$. Let $S'_i$ be the
boundary of the core of $C_i$. It is the central two-sphere of an
$\epsilon$-neck $N'_i$ in $C_i$. We orient the $s$-direction of
$N'_i$ so that the non-compact end of $C_i$ lies off the positive
end of $N'_i$. We denote by $h_{i-1}$ the scale of $N_{i-1}$ and by
$h_i'$ the scale of $N'_i$. By Lemma~\ref{lem:epsilon-neck-curvature-pinching} the ratio
$h_{i-1}/h'_i$ is between $0.99$ and $1.01$. Suppose that $S'_i$ is
disjoint from $C_{i-1}$. Then one of the complementary components of
$S'_i$ in $M$ contains $C_{i-i}$, and of course, one of the
complementary components of $S'_i$ is the core of $C_i$. These
complementary components must be the same, for otherwise $C_{i-1}$
would be disjoint from the core of $C_i$ and hence the intersection
of $C_{i-1}$ and $C_i$ would be contained in $N_i$. This cannot
happen since  $C_{i-1}$ is contained in $C_i$. Thus, if $S'_i$ is
disjoint from $C_{i-1}$, then the core of $C_i$ contains $C_{i-1}$.
 This means that the distance from $x_0$ to the
complement of $C_i$ is greater than the distance of $x_0$ to the
complement of $C_{i-1}$ by an amount equal to the width of $N_i$.
Since the scale of $N_i$ is at least $C^{-1/2}R(x_0)^{-1/2}$ (see
(5) of Definition~\ref{epscap}), it follows from
Corollary~\ref{cor:neck-distance-estimates} that this width is at least
$2(0.99)C^{-1/2}R(x_0)^{-1/2}\epsilon^{-1}$.

Next suppose that $S'_i$ is contained in $C_{i-1}$. Then one of the
complementary components $A$ of $S'_i$ in $M$ has closure contained
in $C_{i-1}$. This component cannot be the core of $C_i$ since the
closure of the core of $C_i$ contains a point of the frontier of
$C_{i-1}$ in $M$. Thus, $A$ contains $N_i$. Of course, $A\not= N_i$
since the frontier of $A$ in $M$ is $S'_i$ whereas $N_i$ has two
components to its frontier in $M$. This means that $C_i$ does not
contain $A$, which is a contradiction since $C_i$ contains $C_{i-1}$
and $A\subset C_{i-1}$.

Lastly, we suppose that $S'_i$ is neither contained in $C_{i-1}$ nor
in its complement. Then $S'_i$ must meet $N_{i-1}$.  According to
Proposition~\ref{prop:overlapping-necks-properties} the $s$-directions in $N_{i-1}$ and
$N'_i$ either almost agree or are almost opposite. Let $x\in
S'_i\cap \partial N_{i-1}$ so that $s_{N'_i}(x)=0$. Move from $x$
along the $s_{N'_i}$-direction that moves into $N_{i-1}$ to a point
$x'$ with $|s_{N_i}(x')|=(0.05)\epsilon^{-1}$. According to
Proposition~\ref{prop:overlapping-necks-properties}
$(0.94)\epsilon^{-1}<s_{N_{i-1}}(x')<(0.96)\epsilon^{-1}$.
 Let $S'(x')$ be the two-sphere in the neck structure for $N'_i$
 through this point. According to Proposition~\ref{prop:overlapping-necks-properties},
 $S'(x')\subset N_{i-1}$, and $S'(x')$ is isotopic in $N_{i-1}$ to its
 central two-sphere. One of the complementary components of $S'(x')$
 in $C_i$, let us call it $A'$, is diffeomorphic to $S^2\times (0,1)$. Also, one of the
 complementary components $A$ of $S'(x')$ in $M$ contains the core of
 $C_{i-1}$. As before, since $C_{i-1}\subset C_i$, the complementary
 component $A$ cannot meet $C_i$ in $A'$. This means that
 the $s_{N_{i-1}}$- and $s_{N'_i}$-directions almost line up along
 $S'(x')$. This means that
 $S'(x')=s_{N'_i}^{-1}(-(0.05)\epsilon^{-1})$. Since the diameter of
 $S'(x')$ is less than $2\pi h_{i-1}$, and since
  $s_{N_{i-1}}(x')\ge (0.94)\epsilon^{-1}$, it follows
 that $S'(x')\subset
 S_{N_{i-1}}^{-1}((0.9\epsilon^{-1},\epsilon^{-1}))$.
Since the distance from $S'_i$ to the central two-sphere is at least
$(0.99)\epsilon^{-1}h'_i$,  It follows from Corollary~\ref{cor:neck-distance-estimates}
that the central two-sphere of $N_i$ is disjoint from $C_{i-1}$ and
lies off the positive end of $N_{i-1}$. This implies that the
distance from $x_0$ to the complement of $C_i$ is greater than the
distance from $x_0$ to the complement of $C_{i-1}$ by an amount
bounded below by the distance from the central two-sphere of $N_i$
to its positive end. According to Corollary~\ref{cor:neck-distance-estimates} this
distance is at least $(0.99)\epsilon^{-1}h_i$, where $h_i$ is the
scale of $N_i$. But we know that  $h_i\ge C^{-1/2}R(x_0)^{-1/2}$.

Thus, all cases either lead to a contradiction or to the conclusion
that the distance from $x_0$ to the complement of $C_i$ is at least
a fixed positive amount (independent of $i$) larger than the
distance from $x_0$ to the complement of $C_{i-1}$.
 Since the diameter of any $(C,\epsilon)$-cap is
uniformly bounded, this contradicts the existence of an infinite
chain $C_0\subset C_1\subset \cdots$ contrary to the claim. This
completes the proof of the claim.
\end{proof}

\begin{claim}[An isotopic separating sphere yields a component of $M$]
\label{claim:isotopic-sphere-gives-component}
\label{sinside}
Suppose that there is a two-sphere $\Sigma\subset N'$ contained in
the closure of the positive half of $N'$ and also contained in
$\widetilde C$. Suppose that $\Sigma$ is isotopic in $N'$ to the
central two-sphere $S'$ of $N'$. Then $\widetilde C\cup C'$ is a
component of $M$, a component containing $X$.
\end{claim}

\begin{proof}
\sketch
$\Sigma$ separates $\widetilde C$ into two components: $A$, which
has compact closure in $\widetilde C$, and $B$, containing the end
of $\widetilde C$. The two-sphere $\Sigma$ also divides $C'$ into
two components. Since $\Sigma$ is isotopic in $N'$ to $S'$, the
complementary component $A'$ of $\Sigma$ in $C'$ with compact
closure contains the closure of the core of $C'$. Of course, the
frontier of $A$ in $M$ and the frontier of $A'$ in $M$ are both
equal to $\Sigma$. If $A=A'$, then the closure of the core of $C'$
is contained in the closure of $A$ and hence is contained in
$\widetilde C$, contradicting our assumption that $C'$ contains a
point of the frontier of $\widetilde C$. Thus, $A$ and $A'$ lie on
opposite sides of their common frontier. This means that $\overline
A\cup \overline A'$ is a component of $M$. Clearly, this component
is also equal to $\widetilde C\cup C'$. Since $X$ is connected and
this component contains a point $x_0$ of $X$, it contains $X$. This
completes the proof of Claim~\ref{claim:isotopic-sphere-gives-component}.
\end{proof}

\begin{claim}[Adjacent cap union is a component of $M$]
\label{claim:cap-union-is-component}
\label{cccomponent}
If $C'$ is a $(C,\epsilon)$-cap whose core contains a point of the
frontier of $\widetilde C$, then $\widetilde C\cup C'$ is a
component of $M$ containing $X$.
\end{claim}

\begin{proof}
\sketch
\uses{claim:isotopic-sphere-gives-component,prop:overlapping-necks-properties,cor:neck-distance-estimates}
We suppose that $\widetilde C$ is the union of $C_0$ and a balanced
chain $N_0,\ldots,N_k$ of $\epsilon$-necks. We orient this chain so
that $C_0$ lies off the negative end of each of the $N_i$.
 Let $S'$ be the boundary of the core of $C'$ and let $N'$ be an
$\epsilon$-neck contained in $C'$ whose central two-sphere is $S'$.
We orient the direction $s_{N'}$ so that the positive direction
points away from the core of  $C'$. The first step in proving this
claim is to establish the following.


Now we return to the proof of Claim~\ref{claim:cap-union-is-component}. We consider
three cases.



\noindent{\bf First Subcase: $S'\subset \widetilde C$.} In this case
we apply Claim~\ref{claim:isotopic-sphere-gives-component} to see that $\widetilde C\cup C'$ is a
component of $M$ containing $X$.


 \noindent {\bf Second Subcase: $S'$ is disjoint from
$\widetilde C$.} Let $A$ be the complementary component of $S'$ in
$M$ containing $\widetilde C$. The intersection of $A$ with $C'$ is
either the core of $C'$ or is a submanifold of $C'$ diffeomorphic to
$S^2\times (0,1)$. The first case is not possible since it would
imply that the core of $C'$ contains $\widetilde C$ and hence
contains $C_0$, contrary to the way we chose $C_0$. Thus, the core
of $C'$ and the the complementary component $A$ containing
$\widetilde C$ both have $S'$ as their frontier and they lie on
opposite sides of $S'$. Since the closure of the core of $C'$
contains a point of the frontier of $\widetilde C$, it must be the
case that $S'$ also contains a point of this frontier. By
Proposition~\ref{prop:overlapping-necks-properties},  the neck $N'\subset C'$ meets
$N_{k}$ and there is a two-sphere $\Sigma\subset N'\cap N_{k}$
isotopic in $N'$ to $S'$ and isotopic in $N_{k}$ to the central
two-sphere of $N_{k}$. Because $N_{k}\subset \widetilde C$ and
$\widetilde C$ is disjoint from the core of $C'$, we see that
$\Sigma$ is contained in the positive half of $N'$. Applying
Claim~\ref{claim:isotopic-sphere-gives-component} we see that $\widetilde C\cup C'$ is a component
of $M$ containing $X$.


\noindent{\bf Third Subcase: $S'\cap \widetilde C\not=\emptyset$ and
$S'\not\subset \widetilde C$.} Clearly, in this case $S'$ contains a
point of the frontier of $\widetilde C$ in $M$, i.e., a point of the
frontier of the positive end of $N_k$ in $M$. Since $N_k\cap
N'\not=\emptyset$, by Lemma~\ref{lem:epsilon-neck-curvature-pinching} the scales of $N_k$ and
$N'$ are within $1\pm 0.01$ of each other, and hence the diameter of
$S'$ is at most $2\pi$ times the scale of $N_k$. Since the central
two-sphere $S'$ of $N'$ contains a point in the frontier of the
positive end of $N_k$, it follows from Corollary~\ref{cor:neck-distance-estimates} that $S'$
is contained on the positive side  of the central two-sphere of
$N_k$ and that the frontier of the positive end of $N_k$ is
contained in $N'$. By Proposition~\ref{prop:overlapping-necks-properties} there is a
two-sphere $\Sigma$ in the neck structure for $N'$ that is contained
in $N_k$ and is isotopic in $N_k$ to the central two-sphere from
that neck structure. Let $A$ be the complementary component of
$\Sigma$ in $M$ that contains $\widetilde C\setminus N_k$.
 If the complementary component of $\Sigma$ that contains
$C'\setminus N'$ is not $A$, then $\widetilde C\cup C'$ is a
component of $M$ containing $X$. Suppose that $A$ is also the
complementary component of $\Sigma$ in $M$ that contains
$C'\setminus N'$. Of course, $A$ is contained in the core of $C'$.
If $k\ge 1$, we see that $A$ and hence the core of $C'$ contains
$\widetilde C\setminus N_k$, which in turn contains the core of
$C_0$. This contradicts our choice of $C_0$. If $k=0$, then
$C_0=A\cup (N_0\cap (M\setminus A))$. Of course, $A\subset C'$.
Also, the frontier of $N_0\cap(M\setminus A)$ in $M$ is the union of
$A$ and the frontier of the positive end of $N_0$ in $M$. But we
have already established that the frontier of the positive end of
$N_0$ in $M$ is contained in $N'$. Since $A\subset C'$, it follows
that all of $C_0$ is contained in $C'$. On the other hand, there is
a point of the frontier of $C_0$ contained in the closure of the
core of $C'$. This then contradicts our choice of $C_0$.

 This
completes the analysis of all the cases and hence completes the
proof of the adjacent cap union assertion.
\end{proof}


As an immediate corollary we have:


\begin{proposition}[Topological classification of manifolds covered by necks and caps]
\label{prop:classification-covered-by-necks-caps}
\uses{prop:connected-set-in-caps-and-necks, def:epsilon-neck, def:c-epsilon-cap}
\label{epstopology}

For all $\epsilon>0$ sufficiently small the following holds. Suppose
that $(M,g)$ is a connected Riemannian manifold such that every
point is either contained in the core of a $(C,\epsilon)$-cap in $M$
or is the center of an $\epsilon$-neck in $M$. Then one of the
following holds:
\begin{enumerate}
\item[(1)] $M$ is diffeomorphic to $S^3$, $\Ar P^3$ or $\Ar P^3\#\Ar
P^3$,
and $M$ is either a double $C$-capped $\epsilon$-tube or is the
union of two $(C,\epsilon)$-caps.
\item[(2)] $M$ is diffeomorphic to  $\Ar^3$ or $\Ar P^3\setminus \{\mathrm{point}\}$,
and $M$ is either a $(C,\epsilon$-cap or a $C$-capped
$\epsilon$-tube.
\item[(3)] $M$ is diffeomorphic to $S^2\times \Ar$ and is an
$\epsilon$-tube.
\item[(4)] $M$ is diffeomorphic to an $S^2$-bundle over $S^1$ and is an $\epsilon$-fibration.
\end{enumerate}
\end{proposition}
